\documentclass[12pt, a4paper, faculty=we]{ugent-doc}

\usepackage[dutch]{babel}

\usepackage{amsfonts}
\usepackage{amssymb}
\usepackage{ dsfont }

\usepackage[T1]{fontenc}
\usepackage[utf8]{inputenc}
% Comment/remove the two lines below to use the default Computer Modern font
\usepackage{libertine}
\usepackage{libertinust1math}

% Mathematics
%-------------

\usepackage{amsmath}
\usepackage{amsthm}


% Figures
%---------
\usepackage{graphicx}
\graphicspath{{./figures/}}

% Bibliography settings
%-----------------------
\usepackage[backend=biber, style=apa, sorting=nyt, hyperref=true]{biblatex}
\addbibresource{./references.bib}

\usepackage{csquotes} % Suggested when using babel+biblatex
% local run procedure: (1) pdfLatex on main, (2) biber on main, (3) pdfLatex on main, (4) pdfLatex on main

% Hyperreferences
%-----------------
\usepackage[colorlinks=true, allcolors=ugentblue]{hyperref}

% Whitespace between paragraphs and no indentation
%--------------------------------------------------
\usepackage[parfill]{parskip} 

\newcommand{\E}[1]{\mathds{E}[#1]}
\newcommand{\h}[1]{\hat{#1}}
\newcommand{\bfZ}{\textbf{Z}}

\newtheorem{stelling}{Stelling}
\newtheorem*{remark}{Opmerking}




%TODO Streamline all notation
%TODO G-Computation deel afwerken
%TODO Lineaire regressie deel afwerken
%TODO Logistische regressie deel afwerken
%TODO Vervang text in align* environment door intertext





\thetitle{\color{ugentblue} ATE evaluation for low sample size using permutation tests with covariate adjustment}

\thesubtitle{Simulation study}

\infoboxa{\bfseries\large Bachelorproef}

% The second infobox at bottom of titlepage
\infoboxb{Promotor en begeleider: 
\newline
\begin{tabular}[t]{l}
    Prof.\ Dr.\ Stijn Vansteelandt\\
    Achille Demares
\end{tabular}
}

\infoboxc{Student: 
\begin{tabular}[t]{lll}
    Ilian Verlee & 02210173 & ilian.verlee@ugent.be
\end{tabular}
}

% The last (bottom) infobox at bottom of titlepage
\infoboxd{Academic year: 2025--2026} % note dash, not hyphen





\begin{document}


\maketitle

\tableofcontents

\newpage
\section{Inleiding}
Hier komt de tekst van de inleiding.

\newpage
\section{Theorie}

\subsection{Waarom permutatie testen?}

\begin{flushleft}
    Hier duiden we onze keuze om te kijken naar permutatie testen. Bij asymptotische theorie wordt er vaak een groot genoege 
    steekproefgrootte veronderstelt zodat men de distributie kan bepalen. Indien dit niet het geval is moet er al geweten zijn welke distributie de test statistiek volgt.
     Bij echte klinische studies is dit echter niet altijd mogelijk, bijvoorbeeld bij studies waarbij maar een aantal patiënten de onderzochte ziekte hebben. 
    Dit leidt vaak tot een reductie in power. Een mogelijke oplossing voor dit probleem is de zogeheten permutatie test.
    Deze heeft geen assumpties op de distributie. De nulhypothese voor de permutatietest is dat alle samples van dezelfde distributie
    komen. We permuteren dan onze data om te zien of dit correct is. Randomisatie en semi-permutatie testen zijn gelijkaardig, maar hebben wel hun eigen nuance.
\end{flushleft}

\subsection{Permutatie vs randomisatie testen}
\begin{flushleft}
    In de literatuur worden de termen randomisatie test en permutatie test vaak door elkaar gebruikt. Er is echter een verschil
tussen de twee, namelijk dat bij een randomisatie test een nieuwe behandelingsvector wordt aangemaakt met een 50\% kans om de behandeling te krijgen,
terwijl bij een permutatie test de behandelingsvector wordt herverdeeld. Een uitgebreidere uitleg en historische context
kunt u terugvinden in \textcite{onghena2018randomization}. We gaan hieronder wat verdere theorie uitwerken over deze twee testen en een derde test, namelijk de semi-permutatie test. 
\end{flushleft}

\subsection{Permutatie testen}

\begin{flushleft}
    We werken enkel de theorie uit voor een two-sample setting, aangezien deze het meest relevant is voor ons onderzoek.
\end{flushleft}

\subsubsection{Opstelling}

\begin{flushleft}
    Neem aan dat we een dataset hebben met \(N\) observaties,waarbij de behandeling willekeurig wordt toegewezen.
    Stel dat willeukerig \(n\) observaties de behandeling krijgen en \(m : = N-n\) observaties de controle groep vormen.
    Zij \(Y^1_i\) de observatie van individu \(i\) in de behandelingsgroep en \(Y^0_i\) de observatie van individu \(i\) in de controle groep.
    Neem aan:
    \[
    \begin{aligned}
        Y^1_1, \ldots, Y^1_n &\quad \text{iid} \quad F_{Y^1}(x), \\
        Y^0_1, \ldots, Y^0_m &\quad \text{iid} \quad F_{Y^0}(x).
    \end{aligned}
    \] met \(H_0 : F_{Y^1}(x) = F_{Y^0}(x)\).
    
\end{flushleft}


\subsubsection{Procedure}

\begin{flushleft}
    We kunnen dan de testprocedure (voor een eenzijdige test) als volgt beschrijven:
    \begin{enumerate}
        \item Wijs willekeurig \(n\) observaties toe aan de behandelingsgroep en \(m : = N-n\) observaties aan de controle groep, waarbij $N$ de grote is van de steekproef.
        \item Voer het experiment uit en bereken een statistiek \(T_{(0)}\) op basis van de verkregen data en een statistiek \(T\).
        \item Bekijk \(\mathbf{Y} := (Y^1_1, \ldots, Y^1_n, Y^0_1, \ldots, Y^0_m)\) als een vector van \(N\) observaties. Sorteer de \(M_n := \binom{N}{n}\) 
        waarden bekomen voor de statistiek \(T\) voor elke gepermuteerde vector \(\mathbf{Y}^*\). Noem deze gesorteerde waarden 
        \(T_{(1)}  \leq \ldots \leq T_{(M_n)}\).
        \item Verwerp \(H_0\) als \(T_{(0)} > T_{(k)}\) met \(k:= M_n-[M_n\cdot\alpha]\). Dit komt overeen met zeggen dat we \(H_0\) verwerpen als 
        de bekomen p-waarde (hieronder gedefinieerd) kleiner is dan \(\alpha\).
    \end{enumerate}
    De p-waarde van deze test is gegeven door:
    \[
        p = \frac{1}{M_N} \sum_{i=1}^{M_n} I(T_{(i)} \geq T_{(0)})
    \]
    waarbij \(I\) de indicatorfunctie is.
\end{flushleft}

\begin{flushleft}
    Om hiervan een tweezijdige test te maken vervang je de \(T_{(i)}\)'s door hun absolute waarde \(|T_{(i)}|\) in de bovenstaande procedure.
\end{flushleft}

\subsubsection{Geldigheid}

\begin{flushleft}
    Stel \(\mathbf{Z}\) een vector met gezamenlijke verdelingsfunctie \(F_{\mathbf{Z}}\) en steekproefruimte \(S\).
    Zij \(G\) een groep van \(M_n\) transformaties van $S$ naar \(S\) zodat onder de nulhypothese de verdeling van elke \(g_i(\mathbf{Z})\) \((g_i\in G)\) precies dezelfde is als de verdeling van \(\mathbf{Z}\).

    Zij \(T(z)\) een reëelwaardige functie op S zodat voor elke \(z\in S\):
    \[T_{(1)}(z)\leq T_{(2)}(z) \leq \dots \leq T_{(M_n)}(z)\]
    de geordende waarden zijn van \(T(g_i(z))\). Gegeven \(\alpha\), definieer:
    \[k:= M_n-\bigl \lfloor M_n\cdot\alpha \bigr \rfloor\] met \(\lfloor \cdot \rfloor\) de floor functie.

    Stel dat \(M_n^+(z)\) en \(M_n^0(z)\) de getallen zijn van \(T_{(j)}(z)\) ,die groter zijn dan \(T_{(k)}(z)\) en gelijk zijn aan \(T_{(k)}(z)\), respectievelijk. Definieer:
    \[a(z) = \frac{M_n\alpha-M_n^+(z)}{M_n^0(z)}\]
    Definieer vervolgens:
    \[\phi(z)= \begin{cases}
        1,\hspace{1cm}T(z)> T_{(k)}(z)\\
        a(z),\hspace{0.50cm}T(z)=T_{(k)}(z)\\
        0,\hspace{1cm}T(z) < T_{(k)}(z)
    \end{cases}\]
    \begin{stelling}[Hoeffding]
        \label{theorem: Hoeffding}
        Stel de data \(\bfZ= (Z_1,\dots, Z_N) \) en de groep \(G\) van transformaties zodanig dat \(g_i(\bfZ) \overset{d}{=} \bfZ\) voor elke \(g_i \in G\) onder \(H_0\) . Dan heeft de test
    gedefinieerd door \(\phi(\bfZ)\) grootte \(\alpha\), m.a.w. \(P(reject \, | \, H_0) = \alpha\).
    \end{stelling}
    \begin{proof}
        Merk eerst op dat door de definitie van $a(z)$ en $\phi$ dat voor elke $z \in S$:
        \[\frac{1}{M_n}\sum_{i=1}^{M_n}\phi(g_i(z))=\frac{M_n^+(z)+a(z)M_n^0(z)}{M_n}=\alpha\]

        Omdat \(g_i(\bfZ) \overset{d}{=} \bfZ\) en \(G\) is een groep, \(\E{\phi(\bfZ) \, | \, H_0}=\E{\phi(g_i(\bfZ)) \, | \, H_0}\) voor elke \(i\)
        \begin{align*}
            P(verwerpen \, | \, H_0)=\E{\phi(\bfZ) \, | \, H_0}&=\frac{1}{M_n}\sum_{i=1}^{M_n}\E{\phi(g_i(\bfZ)) \, | \, H_0}\\
            &=\E\left[ \frac{1}{M_n}\sum_{i=1}^{M_n}\phi(g_i(z)) \, | \, H_0 \right]=\alpha
        \end{align*}
    \end{proof}
\end{flushleft}




\subsection{Randomisatie testen}

\subsubsection{Intuïtie}
\begin{flushleft}
    Een randomisatie test is een test waarbij een nieuwe behandelingsvector wordt gegenereerd waarbij de kans de echte kans is om
    de behandeling te krijgen, onder de aanname dat de behandeling willekeurig is verdeeld, namelijk 1/2. 
\end{flushleft}

\subsubsection{Procedure}

\begin{flushleft}
    We kunnen dan de testprocedure (voor een eenzijdige test), gelijkaardig aan de vorige testen als volgt uitwerken:
    \begin{enumerate}
        \item Wijs willekeurig \(n\) observaties toe aan de behandelingsgroep en \(m : = N-n\) observaties aan de controle groep, waarbij $N$ de grote is van de steekproef.
        \item Voer het experiment uit en bereken een statistiek \(T_{(0)}\) op basis van de verkregen data en een statistiek \(T\).
        \item Bekijk \(\mathbf{Y} := (Y^1_1, \ldots, Y^1_n, Y^0_1, \ldots, Y^0_m)\) als een vector van \(N\) observaties. Sorteer de \(M_n := 2^N\) 
        waarden bekomen voor de statistiek \(T\) voor elke nieuwe vector \(\mathbf{Y}^* \sim  Bernoulli \, ( \, 1/2 \, )\).
        Noem deze gesorteerde waarden 
        \(T_{(1)}  \leq \ldots \leq T_{(M_n)}\).
        \item Verwerp \(H_0\) als \(T_{(0)} > T_{(k)}\) met \(k:= M_n-[M_n\cdot\alpha]\). Dit komt overeen met zeggen dat we \(H_0\) verwerpen als 
        de bekomen p-waarde (zoals eerder gedefinieerd) kleiner is dan \(\alpha\).
    \end{enumerate}
\end{flushleft}

\begin{flushleft}
    Om hier terug een tweezijdige test van te maken vervang je de \(T_{(i)}\)'s door hun absolute waarde \(|T_{(i)}|\) in de bovenstaande procedure.
\end{flushleft}

\subsubsection{Geldigheid}


\begin{flushleft}
Zij $N \in \mathds{N}$ en:
\begin{itemize}
    \item $\mathbf{Y}^1, \, \mathbf{Y}^0, \, \mathbf{X}  \in \mathds{R}^N $ vectoren
    \item $\mathbf{Z}$ een vector met componenten $Z_i \in \{0,1\}^N $, onafhankelijk van $\mathbf{Y}^1$ en $\mathbf{Y}^0$
    \item Een vector $\mathbf{Y} = \mathbf{Z} \times \mathbf{Y}^1 + (1-\mathbf{Z}) \times \mathbf{Y}^0$
\end{itemize}

Zij $T(\bfZ,\mathbf{Y})$ een reëelwaardige scalaire functie, definieer dan:
\begin{equation}
    \label{eq: pRand}
    p(T,Z,Y) = \sum_{\mathbf{Z}'} F_{\mathbf{Z}}(\mathbf{Z}')\cdot I(T(\mathbf{Z}' , \mathbf{Y}) \geq T(\mathbf{Z} , \mathbf{Y}))
\end{equation}
waarbij $\mathbf{Z}' \overset{d}{=} \bfZ$
\end{flushleft}

\begin{stelling} \footnote{\cite{gherardi2023}}
    \label{theorem: Fisher}
    Als $\mathbf{Y}^1 = \mathbf{Y}^0$ (m.a.w. onder $H_0$), dan:
    \[P(p(T,Z,Y)\leq \alpha) \leq \alpha \]
\end{stelling}
\begin{proof}
    Stel $\mathbf{Z}' \overset{d}{=} \bfZ$ en definieer $\mathbf{Y}' = \mathbf{Z}' \times \mathbf{Y}^1 + (1-\mathbf{Z}') \times \mathbf{Y}^0$. Gegeven $T_0 \in \mathds{R}$
    , berekenen we:
    \begin{align*}
        P(T(\mathbf{Z}', \mathbf{Y}') \geq T_0 \, | \mathbf{Y}^0, \mathbf{Y}^1) &= \E{I(T(\mathbf{Z}', \mathbf{Y}') \geq T_0) \, | \mathbf{Y}^0, \mathbf{Y}^1} \\
        &= \sum_{\mathbf{Z}'} F_{\mathbf{Z}}(\mathbf{Z}')\cdot I(T(\mathbf{Z}' , \mathbf{Y}') \geq T_0)
    \end{align*}
    waar we voor de eerste gelijkheid gebruik maken van $\E{I_A} = P(A)$ voor een verzameling $A$. Bij de tweede overgang gebruiken we dat $\mathbf{Z}$ discreet is aangezien
    er maar $2^N$ keuzes zijn voor deze vector. Als nu de nulhypothese geldt, dus  $\mathbf{Y}^1 = \mathbf{Y}^0$, dan hebben we $T(\mathbf{Z}', \mathbf{Y}') = 
    T(\mathbf{Z}', \mathbf{Y})$, zodat we $\mathbf{Y}'$ door $\mathbf{Y}$ kunnen vervangen hierboven. Als we bovendien $T_0 = T(\bfZ, \mathbf{Y})$ kiezen dan wordt het
    rechterlid $p(T,Z,Y)$. We krijgen dus:
    $$p(T,Z,Y)=P(T(\mathbf{Z}', \mathbf{Y}') \geq T(\bfZ, \mathbf{Y}) \, | \mathbf{Y}^0, \mathbf{Y}^1)$$
    Dus:
    \begin{align*}
        P(p(T,Z,Y)\leq \alpha | \mathbf{Y}^0, \mathbf{Y}^1)  &= P(P(T(\mathbf{Z}', \mathbf{Y}') \geq T(\bfZ, \mathbf{Y}) )\leq \alpha | \mathbf{Y}^0, \mathbf{Y}^1)\\
                                &= P(T(\mathbf{Z}', \mathbf{Y}') \leq F^{-1}_{T(\bfZ, \mathbf{Y})} (\alpha) | \mathbf{Y}^0, \mathbf{Y}^1)\\
                                &= F_{T(\mathbf{Z}', \mathbf{Y}') | \mathbf{Y}^0, \mathbf{Y}^1}(F^{-1}_{T(\bfZ, \mathbf{Y})} (\alpha))\\
    \text{We weten dat $T(\mathbf{Z}', \mathbf{Y}') | \mathbf{Y}^0, \mathbf{Y}^1 \overset{d}{=} T(\bfZ, \mathbf{Y})$, dus:}\\
                                &= F_{T(\bfZ, \mathbf{Y})}(F^{-1}_{T(\bfZ, \mathbf{Y})} (\alpha))\\
                                &\leq \alpha
    \end{align*}
    Aangezien dit geldt voor elke $\mathbf{Y}^0, \mathbf{Y}^1$ geldt het gestelde.
\end{proof}

\begin{flushleft}
    Nemen we aan dat $\bfZ \sim Bernoulli(p)$ verdeeld is dan krijgen we: 
    $$F_{\mathbf{Z}}(\mathbf{z}) = \prod_{i=1}^{N} p^{z_i} (1-p)^{(1-z_i)}  $$
    Als we aannemen dat $p = 1/2$, dan krijgen we:
    $$F_{\mathbf{Z}}(\mathbf{z}) = \frac{1}{2^N}$$
    Vullen we dit in in (\ref{eq: pRand}), krijgen we:
    $$p(T,Z,Y) = \frac{1}{2^N} \sum_{\mathbf{Z}'} I(T(\mathbf{Z}' , \mathbf{Y}) \geq T(\mathbf{Z} , \mathbf{Y}))$$
    We kunnen nu zien dat dit exact hetzelfde is als de p-waarde voor de randomisatie test eerder gedefinieerd.
\end{flushleft}


\subsection{Semi-permutatie testen}

\subsubsection{Intuïtie}
\begin{flushleft}
    Een semi-permutatie test is een test waarbij er een nieuwe behandelingsvector
    wordt gegeneerd waarbij de kans om de behandeling te krijgen gelijk is aan
    het gemiddelde van de oorspronkelijke behandelingsvector. Deze test is een soort van tussenstap 
    tussen een randomisatie test en een permutatie test.
\end{flushleft}

\subsubsection{Procedure}

\begin{flushleft}
    We kunnen dan de testprocedure (voor een eenzijdige test), gelijkaardig aan de permutatie test omschrijven als volgt:
    \begin{enumerate}
        \item Wijs willekeurig \(n\) observaties toe aan de behandelingsgroep en \(m : = N-n\) observaties aan de controle groep, waarbij $N$ de grote is van de steekproef.
        \item Voer het experiment uit en bereken een statistiek \(T_{(0)}\) op basis van de verkregen data en een statistiek \(T\).
        \item Bekijk \(\mathbf{Y} := (Y^1_1, \ldots, Y^1_n, Y^0_1, \ldots, Y^0_m)\) als een vector van \(N\) observaties. Sorteer de \(M_n := 2^N\) 
        waarden bekomen voor de statistiek \(T\) voor elke nieuwe vector \(\mathbf{Y}^* \sim  Bernoulli \, ( \, \overline{A} \, )\), waarbij $\overline{A} $ het steekproefgemiddelde is van de treatment vector $\mathbf{A}$. 
        Noem deze gesorteerde waarden 
        \(T_{(1)}  \leq \ldots \leq T_{(M_n)}\).
        \item Verwerp \(H_0\) als \(T_{(0)} > T_{(k)}\) met \(k:= M_n-[M_n\cdot\alpha]\). Dit komt overeen met zeggen dat we \(H_0\) verwerpen als 
        de bekomen p-waarde (zoals bij de permutatie test) kleiner is dan \(\alpha\).
    \end{enumerate}
\end{flushleft}

\begin{flushleft}
    Om hier terug een tweezijdige test van te maken vervang je de \(T_{(i)}\)'s door hun absolute waarde \(|T_{(i)}|\) in de bovenstaande procedure.
\end{flushleft}

\subsubsection{Geldigheid}

\begin{flushleft}
    Deze test voldoet niet aan stelling \ref{theorem: Hoeffding} en stelling \ref{theorem: Fisher}, omdat deze nieuwe vector niet dezelfde
    distributie volgt als onze originele toewijzing.
\end{flushleft}

\subsection{G-Computation}

\subsubsection{ATE}
\begin{flushleft}
    We willen een schatter bekomen voor het average treatment effect (ATE). Dit helpt ons het gemiddelde effect van de behandeling
    op de populatie te onderzoeken. De definitie van de ATE wordt gegeven door:
    \begin{equation}
        \label{eq:orATE}
        \E{Y^1 - Y^0}
    \end{equation}
    waarbij $Y^1$ de groep is die de behandeling heeft gekregen en $Y^0$ de controlegroep.
    \newline
    Nu willen we een schatter voor deze waarde. Hiervoor hebben we de volgende assumpties nodig:
    \begin{enumerate}
        \item Conditional exchangeability:\\
        Dit zegt dat de potentiële uitkomsten onafhankelijk zijn van de toewijzing van de behandeling, geconditioneerd op een
        verzameling baseline covariaten $(L)$. Expliciet kunnen we dit schrijven als: $Y^1, Y^0 \perp A | L$ waarbij $A 
        \in \{0,1\}$ onze toewijzing tot de behandelingsgroep $(1)$ of tot de controlegroep $(0)$ is.

        \item Positiviteit: \\
        $$0 < P(A_I = 1 | L_i) < 1$$
        
        \item Consistentie: \\
        Alle deelnemers aan de studie moeten dezelfde behandeling krijgen, onder dezelfde condities en met dezelfde dosis.

        \item Non-interference: \\
        De uitkomst van een individu wordt niet beïnvloed door de toegewezen behandeling van andere individuen. Kort geschreven:
        $$Y_i(\mathbf{L}) = Y_i(L_i)  $$
    \end{enumerate}
\end{flushleft}

We herschrijven nu (\ref{eq:orATE}) met de regel van herhaalde verwachtingswaarde:

    \begin{align*}
        \E{Y^1 - Y^0} &= \E{\E{Y^1 \, | L} - \E{Y^0 \, | L} }\\
        \intertext{Uit $Y^1, Y^0 \perp A | L$ volgt dan:}\\
                          &= \E{\E{Y^1 \, | A=1, \, L} - \E{Y^0 \, | A=0, \, L} }\\
                          &= \E{\E{Y \, | A=1, \, L} - \E{Y \, | A=0, \, L} }\\
        \intertext{Vervangen we nu de eerste verwachtingswaarde door de sample mean krijgen we:}\\
        ATE        &= \frac{1}{n} \sum_{i=1}^{n} \left( \E{Y_i | A_i = 1, \mathbf{L_i}} - \E{Y_i | A_i = 0, \mathbf{L_i}} \right)\\
    \end{align*}

\begin{flushleft}
    Natuurlijk weten we de echte waarde van $\E{Y \, | A=a, \, L}$ niet, dus gaan we hiervoor een model opstellen om dit te schatten.
    We zullen hiervoor de volgende definitie gebruiken: $$\mu_a := \E{Y | A = a, \mathbf{L}}$$. 

    \subsubsection{Lineair model}
    We kunnen onderstellen dat \(\mu_a\) voldoet aan het volgende lineaire model:
    \[
{\mu}_a = \beta_0 + \beta_1 A + \mathbf{\beta_2}^{\top} \cdot \mathbf{L} + A \, (\mathbf{\beta_3}^{\top} \cdot \mathbf{L})    \]
    Hieruit volgt dat onze ATE schatter gegeven wordt door:
    \begin{equation}
        ATE = \frac{1}{n} \sum_{i=1}^{n} \left( \mu_{1,i} - \mu_{0,i} \right) = \beta_1 + (\mathbf{\beta_3}^{\top} \cdot \mathbf{\overline{L}}_n)
    \end{equation}
    met \(\mu_{a,i} = \beta_0 + \beta_1 a + \mathbf{\beta_2}^{\top} \cdot \mathbf{L_i} + a \, (\mathbf{\beta_3}^{\top} \cdot \mathbf{L_i})\) en \(\mathbf{\overline{L}}_n = \frac{1}{n} \sum_{i=1}^{n} \mathbf{L_i}\).

    
\end{flushleft}

\begin{flushleft}
    We kunnen dan de ATE schatten a.d.h.v. de OLS schatter \(\hat{\beta_1}\) en \(\hat{\beta_3}\) van \(\beta_1\) en \(\beta_3\) respectievelijk.
     We kunnen onze ATE schatter dus als volgt schrijven:
        \[\hat{ATE} = \hat{\beta_1} + (\hat{\mathbf{\beta_3}}^{\top} \cdot \mathbf{\overline{L}}_n)\]
    \begin{remark}
    Als er geen interactie is tussen de behandeling en de covariaten, dan is \(\hat{ATE} = \hat{\beta_1}\).
    We weten dat de OLS schatter een onvertekende schatter is van \(\beta_1\) onder de aanname dat het model correct is gespecificeerd of wanneer de 
    verwachtingswaarde van de behandelingsvariabele lineair is in de andere covariaten .
    Dit impliceert dat de OLS schatter onvertekend is voor de \textbf{geschatte ATE} onder de hierboven genoemde voorwaarden. 
    Hieruit volgt:
    \[
        \E{\hat{ATE}} = \E{\hat{\beta_1}} = \beta_1
    \]
    We weten ook het volgende over G-Computation als het model correct is:
    \[
        \E{\hat{ATE}} = \E{\frac{1}{n} \sum_{i=1}^{n} \left( \mu_{1,i} - \mu_{0,i} \right)} = \E{\E{Y | A = 1, \mathbf{L}} - \E{Y | A = 0, \mathbf{L}}} = ATE \]

    Dus onder de aanname dat het lineaire model correct is gespecificeerd en er geen interactie is tussen de behandeling en de covariaten, is de OLS schatter van \(\beta_1\) een onvertekende schatter van de ATE.
    \end{remark}
\end{flushleft}

\subsubsection{Logistische regressie}
\paragraph[Wat is logistische regressie?]{Logistische regressie is}
Logistische regressie is


\begin{flushleft}
    
\end{flushleft}

\subsubsection{}

\begin{flushleft}
    Voor een logistische regressie gaat het moeilijker zijn om een expliciete vorm
    te vinden voor de ATE schatter. Dit komt doordat de covariaten niet wegvallen in onze uitdrukking
    zoals hieronder gegeven.
\end{flushleft}
\begin{flushleft}
    Zij ons logistisch regressie model gegeven door:
    \[
        \hat{\mu_a} = \frac{1}{1 + e^{-(\beta_0 + \beta_1 A + \vec{\beta_2}^T \vec{L})}}
    \]
    Dan is onze ATE schatter gegeven door:
    \begin{equation}
        \hat{ATE} = \frac{1}{n} \sum_{i=1}^{n} \left( \hat{\mu}_{1,i} - \hat{\mu}_{0,i} \right)
    \end{equation}
    met \(\hat{\mu}_{a,i} = \frac{1}{1 + e^{-(\beta_0 + \beta_1 a + \vec{\beta_2}^T \vec{L_i})}}\).
    \newline \hfill \break
    Als we deze uitdrukking proberen te vereenvoudigen, zien we dat de covariaten niet wegvallen.
    Dit impliceert dat we geen gemakkelijk berekenbare vorm kunnen vinden voor onze ATE schatter.
    We gaan dus a.d.h.v. een predictie van ons logistisch regressie model de ATE schatter berekenen.
\end{flushleft}


\subsection{LASSO}

\begin{flushleft}
    De least absolute schrinkage and selection operator, meer gekend als LASSO, is een manier om variable selection
    uit te voeren. M.a.w. een manier om covariaten te selecteren voor ons model.
\end{flushleft}

\begin{flushleft}
    %Hier komt een sectie over LASSO en over het feit dat een model bekomen door LASSO weldegelijk voldoet aan de gegeven voorwaarden
    %ivm de functie
\end{flushleft}
% Dit moet in een andere sectie komn




\section{Simulaties}
\begin{flushleft}
    In ons onderzoek hebben we gekozen om 5 verschillende testen te vergelijken met elkaar voor het lineaire geval: randomisatie, permutatie, semi-permutatie testen met covariate adjustment, de klassieke t-test en covariate adjustment op de klassieke manier uitgevoerd.
    Hierbij werden de randomisatie, permutatie en semi-permutatie testen ook opgesplitst in twee procedures. Bij de eerste procedure wordt er een model voor elke permutatie gefit a.d.h.v. LASSO en vergelijken we de hieruit bekomen ATE. Bij de tweede procedure berekenen
    we de covariaten die we gaan gebruiken in ons model a.d.h.v. lasso maar fitten het model opnieuw met de covariaten bekomen door de LASSO. Deze procedure wordt ook wel post-LASSO genoemd.
     We vergeleken deze testen dan op basis van hun power en type I error rate.
\end{flushleft}



% Voor bibliography moet er een textcite zijn in de tekst, anders wordt deze niet opgenomen in de bibliography
\nocite{*}
\printbibliography[heading=bibintoc, title={Bibliografie}]

\end{document}